% =====================================================================
%  LAYER B -- The compiled evaluator.
%  Section 6: every code contract in Tier 1, gathered.
%
%  Source: tier1.md 1.1.4 (302-329) + the prefactor/ye_weighted snippets
%          of 1.1.2-1.1.3 (229-239, 286-300); II.2's vflag_sets extraction
%          (965-976); II.4 (1028-1058); II.5 (1060-1087); II.6b
%          (1116-1148); III.3's pf spline + pf_matrix (1429-1472);
%          III.8 db_reverses.py (1714-1745); IV.5/IV.6 vectorisation and
%          screen_map (2183-2185, 2285-2296); V.5's compile/engine
%          snippets (2753-2773); V.6's ledger assertion (2854-2870);
%          V.7's duplicate-link resolution (2917-2928); V.9 the code path
%          (2945-2962).
%  Self-check: the items relocated from Part I (item 4), Part II
%          (items 4-6) and Part V (items 5, 10).
% =====================================================================

\section{The compiled evaluator}
\label{part:evaluator}

Everything so far has been physics. This section is the other half of the
DERIVE\,/\,READ loop: the code contract that turns
\S\ref{part:ratetheory}--\S\ref{part:weak} into dense array algebra over
$10^5$--$10^6$ plasma states. The source document interleaves these fragments
with their derivations; gathering them makes visible what is otherwise easy to
miss --- that the whole \code{fluxes/} package is one design idea applied
repeatedly.

\begin{center}
\begin{tabularx}{\textwidth}{@{}L{4.2cm} L{2.0cm} Y@{}}
\toprule
\textbf{What is implemented} & \textbf{Derived in} & \textbf{Here} \\
\midrule
$R_j$ from $\lambda_j$, arity, $\rho$ powers &
  \S\ref{sec:molar} & \S\ref{sec:code-contract} \\
The seven-coefficient sum &
  \S\ref{sec:sevencoeff} & \S\ref{sec:coefftensor} \\
The v flag as a boolean array &
  \S\ref{sec:labelprops} & \S\ref{sec:vflag-extraction} \\
Post-multiplied pf and screening corrections &
  \S\ref{sec:sevencoeff} & \S\ref{sec:lambda-to-R} \\
$\lambda$'s dynamic range &
  \S\ref{sec:multisets} & \S\ref{sec:lambda-range} \\
$\prod G_{\Rset}/\prod G_{\Pset}$ as a signed incidence matmul &
  \S\ref{sec:pf-in-ratio} & \S\ref{sec:pf-spline} \\
Replacing every v-flag reverse with a \code{DerivedRate} &
  \S\ref{sec:vflag} & \S\ref{sec:db-replacement} \\
$h$ per reaction from deduplicated pairs &
  \S\ref{sec:screening-pairs} & \S\ref{sec:screening-code} \\
Bilinear interpolation of the weak tables &
  \S\ref{sec:weaktable} & \S\ref{sec:tabular-path} \\
The lepton ledger as an assertion &
  \S\ref{sec:ledger} & \S\ref{sec:ledger-enforcement} \\
Tabular-wins duplicate resolution &
  \S\ref{sec:tablefamilies} & \S\ref{sec:duplicate-links} \\
\bottomrule
\end{tabularx}
\end{center}

% =====================================================================
\subsection{The code contract}
\label{sec:code-contract}

\code{CompiledNetwork} stores exactly the pieces of \eqref{eq:grossflux}:

\begin{center}
\begin{tabularx}{\textwidth}{@{}L{3.0cm} L{3.4cm} Y@{}}
\toprule
\textbf{Field} & \textbf{Shape} & \textbf{Meaning} \\
\midrule
\code{prefactor} & (n\_rxn,) & $1/\prod c_s!$ \\
\code{dens\_exp} & (n\_rxn,) & n\_reactants $-$ 1 ($+1$ if \Ye-weighted) \\
\code{reactant\_idx} & (n\_rxn, max\_nr) & species indices,
  sentinel $=$ n\_species \\
\code{ye\_weighted} & (n\_rxn,) bool & the \rhoye{} carriers \\
\bottomrule
\end{tabularx}
\end{center}

and \code{engine.evaluate\_fluxes} assembles them:

\begin{srclisting}
R = lam * cn.prefactor[:, None] * rho[None, :] ** cn.dens_exp[:, None]
...
Ypad = np.vstack([Y, np.ones((1, n_states), dtype=np.float64)])  # sentinel row = 1.0
for k in range(cn.reactant_idx.shape[1]):
    R *= Ypad[cn.reactant_idx[:, k]]
\end{srclisting}

The sentinel trick is worth pausing on: reactions have different arities, so the
reactant index array is ragged. Padding with an index that points at a row of
ones turns the ragged product into a fixed-length dense one, at the cost of one
wasted multiply per unused slot.

\begin{result}
This is the recurring shape of the whole \code{fluxes/} package --- \emph{turn
per-reaction control flow into dense array algebra}, which is what makes
$7\times10^5$ states\,/\,min\,/\,core possible \resultsref{2026-07-10}.
\end{result}

The two halves of the \Ye-weighting rule of \S\ref{sec:molar} are mirrored
across compile and evaluate:

\begin{srclisting}
ye_weighted[j] = bool(getattr(rate, "use_ye_weighting", False))   # compile.py:285
...
R[cn.ye_weighted] *= ye[None, :]                                  # engine.py:179
\end{srclisting}

and \code{prefactor} --- the pure combinatorial factor of
\S\ref{sec:double-counting} --- is built once at compile time
(\code{fluxes/compile.py:286}) and never touched again, because it has no
temperature dependence and therefore no business inside the coefficient tensor.

% =====================================================================
\subsection{The coefficient tensor}
\label{sec:coefftensor}

The compile step's job is to turn a ragged object graph (rates, each with a list
of sets, each with seven coefficients) into flat arrays:

\begin{center}
\begin{tabularx}{\textwidth}{@{}L{4.4cm} Y@{}}
\toprule
\textbf{Field} & \textbf{Meaning} \\
\midrule
\code{coeffs} (n\_sets, 7) & every set of every rate, stacked \\
\code{set\_owner} (n\_sets,) & which reaction column owns each set \\
\code{set\_starts}, \code{owned\_cols} & \code{reduceat} boundaries for the
  segment sum \\
\code{pf\_splines}, \code{pf\_matrix} & per-nucleus log-pf splines $+$ the $\pm$
  multiplicity-count incidence matrix ($+n$ per reactant occurrence, $-n$ per
  product occurrence) \\
\code{tab\_*} & per-tabular-rate grids and 2-D log-rate arrays \\
\code{screen\_pairs}, \code{screen\_map} & deduplicated $(Z,A)$ pairs $+$
  per-reaction multiplicity \\
\bottomrule
\end{tabularx}
\end{center}

Two design points worth internalising, because they recur:

\begin{enumerate}
\item \textbf{Deduplicate over pairs, not reactions.} Many reactions share the
      same screening pair (everything capturing an $\alpha$ on a $Z = 14$
      nucleus). \code{screen\_pairs} holds unique pairs; \code{screen\_map} is a
      sparse (n\_rxn $\times$ n\_pairs) multiplicity matrix. Screening is then
      one evaluation per unique pair per state, plus a sparse matmul. Same idea
      for \code{pf\_splines}\,/\,\code{pf\_matrix} over nuclei.
\item \textbf{The pf incidence matrix carries signs, and the signs are of the
      \emph{source} rate.} In \code{compile.py:232-239}, source-rate
      \emph{reactants} get $+1$ and \emph{products} $-1$ --- because the
      correction being applied is the ratio
      $\prod G_{\rm reactants}/\prod G_{\rm products}$ of the \textbf{forward}
      rate, applied to the derived reverse (\eqref{eq:dbratio},
      \S\ref{sec:pf-in-ratio}). Getting that sign backwards would invert every
      partition-function correction while leaving all shapes and magnitudes
      plausible. This is the kind of bug the $\kappa$-at-NSE test
      (\S\ref{sec:oracles-db}) exists to catch.
\end{enumerate}

% =====================================================================
\subsection{The v flag, extracted}
\label{sec:vflag-extraction}

Character 5 of \code{labelprops} (\S\ref{sec:labelprops}) becomes a boolean
array at compile time and is handed straight to the gate:

\begin{srclisting}
vflag_sets = np.array([lp[5] == "v" for lp in set_labelprops], dtype=bool)  # compile.py:250
assert_pf_gate(vflag_sets, context=f"{network} compiled set tensor")
\end{srclisting}

pynucastro's \code{derived\_from\_inverse} boolean surfaces the same fact at the
\code{Rate} level and is exported into the $\nu$ npz, so the gate can be
evaluated either from a live collection or from an artifact on disk
(\S\ref{sec:guard-pf}).

% =====================================================================
\subsection{From \texorpdfstring{$\lambda$}{lambda} to the gross flux}
\label{sec:lambda-to-R}

Equation \eqref{eq:grossflux}, implemented. Two subtleties visible only in the
code:

\paragraph{Screening and pf enter as a post-multiplication.}
\code{evaluate\_lambda} computes the set sum first, then multiplies by
$\exp(\text{corr})$:

\begin{srclisting}
lam[cn.owned_cols] = np.add.reduceat(np.exp(log_sets), cn.set_starts, axis=0)
...
corr = cn.pf_matrix @ log_pf          # DerivedRate columns
corr = corr + cn.screen_map @ log_scor
lam *= np.exp(corr)
\end{srclisting}

This is algebraically identical to adding inside each set's exponent --- provided
the correction is the same for all sets of a column, which it is (both pf and
screening are properties of the \emph{reaction}, not of the fit component). The
module docstring says exactly this, and says it is ``fp-equal to ${\sim}1$
ulp''. That equivalence is precisely what the $5\times10^{-12}$ oracle measures
(\S\ref{sec:oracle-compile}).

\paragraph{Tabular weak rates bypass the coefficient tensor entirely} --- they
are interpolated, then written into the same \code{lam} array
(\S\ref{sec:tabular-path}). One array, two completely different evaluation
mechanisms, unified by column index. That unification is the design decision
that makes everything downstream (\code{stoich.nu @ R}, the pair map, $\kappa$)
indifferent to rate provenance.

\paragraph{Column identity is the invariant that makes it safe.}
\code{compile\_network} asserts the compiled column order equals
\code{stoich.rate\_fnames}, i.e.\ the $\nu$ export order
(\code{compile.py:173-183}). If that ever drifts, every flux is silently
attributed to the wrong reaction --- a failure mode with no numerical signature.
Hence the assertion rather than a comment.

% =====================================================================
\subsection{The numerical range of \texorpdfstring{$\lambda$}{lambda}, and why nothing overflows}
\label{sec:lambda-range}

A deep-dive worth doing once, because it explains several design choices at
once. Measured on \msn{80}, screening off, $\rho = 10^8$ \derivedhere{}:

\begin{center}
\begin{tabular}{@{}lllrr@{}}
\toprule
\textbf{\Tn} & \textbf{$\lambda$ range ($\log_{10}$)} & \textbf{span}
  & \textbf{$\lambda = 0$} & \textbf{$\lambda$ non-finite} \\
\midrule
1.6 & $-63.83 \ldots +11.61$ & \textbf{75 decades} & 0 & 0 \\
3.3 & $-30.61 \ldots +12.93$ & 44 decades & 0 & 0 \\
5.0 & $-24.59 \ldots +14.01$ & 39 decades & 0 & 0 \\
7.9 & $-20.37 \ldots +16.62$ & 37 decades & 0 & 0 \\
\bottomrule
\end{tabular}
\end{center}

And at the level of individual sets, the exponent $a\cdot\text{basis}$ ranges
from \textbf{$-4.4\times10^6$ to $+38$} across the box edges.

Three readings:

\begin{enumerate}
\item \textbf{Nothing ever overflows.} The maximum exponent ($+38$) is nowhere
      near $\ln(\text{DBL\_MAX}) = 709.78$. The evaluator can compute \code{exp}
      unguarded.
\item \textbf{43 of the 1834 set-evaluations across the two box-edge
      temperatures underflow to exactly 0} (\msn{80} has 917 sets; 7 underflow
      at $\Tn = 1.6$ and 36 at $\Tn = 7.9$). This is \emph{correct} --- those
      are rates whose true value is $e^{-4\times10^6}$, i.e.\ zero to any
      meaning of the word. The compiled path relies on IEEE graceful underflow
      rather than masking, which is the right call: a masked branch would cost a
      comparison on every set at every state to reproduce the answer the
      hardware already gives.
\item \textbf{The 75-decade span at $\Tn = 1.6$ is the reason everything
      downstream is log-space or ratio-based.} It is also the reason the
      conservation gate is phrased relative to \emph{gross} flux rather than
      absolutely (Tier~0's $\max(10^{-12}s,\,10^{-13}G)$ form): an absolute
      tolerance is meaningless against a quantity spanning 75 decades. The
      dynamic range narrows as $T$ rises --- the Gamow exponent
      $\tau \propto T^{-1/3}$ compresses everything --- which is why the box's
      cold corner is the numerically hardest, not the hot one.
\end{enumerate}

% =====================================================================
\subsection{The pf spline, and what happens past the table}
\label{sec:pf-spline}

Both \code{fluxes/compile.py} and \code{qse/coeffs.py} rebuild the same object:

\begin{srclisting}
InterpolatedUnivariateSpline(pf.T9_points, pf.log_pf_data,
                             k=pf.interpolant_order, ext="const")
\end{srclisting}

\code{ext="const"} means \textbf{constant extrapolation} --- past the last
tabulated \Tn{} the spline holds its final value rather than continuing the
(steeply rising) trend. Three notes:

\begin{enumerate}
\item Constant extrapolation of $\log G$ is conservative: it
      \emph{under}-estimates $G$ at high $T$ rather than diverging. The
      alternative (\code{ext=0}, spline extrapolation) with a cubic would run
      away.
\item The Rauscher tables extend well past $\Tn = 7.9$, so the box never
      triggers it. Like the \code{smooth\_clip} in the screening module
      (\S\ref{sec:chugunov}), it is a guard that should never fire, present so
      that out-of-box calls degrade rather than explode.
\item \textbf{The two rebuilds must agree.} Rate-side pf
      (\S\ref{sec:pf-in-ratio}) and Saha-side pf (\code{qse/coeffs.py}) are
      constructed from the same fields with the same spline order and the same
      \code{ext} --- deliberately, because $\kappa$-at-NSE compares a rate ratio
      against a composition computed from the equilibrium side. A mismatch there
      would produce a nonzero $\kappa$ floor with no defect in either component.
      The \code{qse/coeffs.py} docstring makes the point explicitly: using the
      same inputs is what makes the solver cross-check ``a genuine independence
      test of the SOLVER, not the data''.
\end{enumerate}

\textbf{In the code}, the correction is applied as a sparse incidence matmul:

\begin{srclisting}
for nuc in rate.source_rate.reactants:   pf_entries[(j, m)] += 1.0
for nuc in rate.source_rate.products:    pf_entries[(j, m)] -= 1.0   # compile.py:232-239
...
corr = cn.pf_matrix @ log_pf ;  lam *= np.exp(corr)                  # engine.py:133,148
\end{srclisting}

Note again (\S\ref{sec:coefftensor}) that the signs are keyed to the
\textbf{source (forward) rate's} reactants and products, matching
$\prod G_{\Rset}/\prod G_{\Pset}$ in \eqref{eq:dbratio}. The same splines, from
the same Rauscher tables, are rebuilt independently in \code{qse/coeffs.py} for
the Saha side --- and the two are cross-checked, because a pf mismatch between
the rate side and the equilibrium side would make $\kappa$-at-NSE nonzero for
reasons that have nothing to do with the rates.

% =====================================================================
\subsection{\texorpdfstring{\code{db\_reverses.py}}{db reverses.py}: the replacement}
\label{sec:db-replacement}

The construction \S\ref{sec:vflag} mandates --- every v-flag reverse swapped for
a \code{DerivedRate(use\_pf=True)} --- is one function:

\begin{srclisting}
def replace_vflag_reverses(rc, table, info=None) -> tuple[list, ReplaceReport]
\end{srclisting}

Positional replacement --- the returned list preserves collection order, so
$\nu$ is untouched. Three details that are each a lesson:

\begin{enumerate}
\item \textbf{Forward lookup is two-tier.} Canonical directed key within the
      collection first, then a lazily-built
      \code{ReacLibLibrary().linking\_nuclei(...)} fallback --- because a v-flag
      reverse can be in the reconciled set while its forward is not. Measured
      outcome: \textbf{all} forward partners were found in-collection (280/280
      \msn{80}, 672/672 \msn{151}) \resultsref{2026-07-10}, so the fallback
      never fired. It exists because the failure would otherwise be silent.

\item \textbf{Multiset identity is asserted, not assumed.}
      \begin{srclisting}
if Counter(map(str, derived.reactants)) != Counter(map(str, rate.reactants)) or ...:
    raise ValueError("... changed the reactant/product multisets -- nu column identity broken")
      \end{srclisting}
      The replacement is the one operation in the pipeline that swaps a live
      object in a column; the assertion pins the only thing that must not change.

\item \textbf{Failures are reported, not swallowed.} \code{ReplaceReport.failed}
      must be empty for the pf gate to pass. A channel whose spin states are
      missing cannot get a \code{DerivedRate}, and the design decision is to
      \emph{stop} and force an explicit routing decision (mesa\_probe24 values,
      or exclude-with-footnote) rather than fall back to the v-flag.
\end{enumerate}

The \code{missing\_pf\_nuclei} field records nuclei with no Rauscher table, where
pynucastro defaults $\log_{\rm pf} = 0$. Measured: light sector only ($d$,
\nuc{He}{3}, \nuc{He}{4}, \nuc{Li}{7}, \nuc{Be}{7,9,10}, \nuc{B}{8},
\nuc{C}{12,13}, \nuc{N}{13-15}) \resultsref{2026-07-10} --- harmless, since
$G \approx 1$ there for exactly the \S0.2.2 reason (few low-lying levels in
light nuclei).

% =====================================================================
\subsection{Screening, vectorised}
\label{sec:screening-code}

\code{fluxes/screening.py} reimplements pynucastro's scalar jitclass over numpy
state vectors; the physics it evaluates is \eqref{eq:plasmastate}--%
\eqref{eq:chugunov}. The per-reaction bookkeeping is the dedup-over-pairs idea
of \S\ref{sec:coefftensor}:

\begin{srclisting}
for n1, n2 in getattr(rate, "screening_pairs", []):
    key = (n1.Z, n1.A, n2.Z, n2.A)  ... screen_map[j, pair_index[key]] += 1
\end{srclisting}

and the per-reaction correction is $\sum$ over its pairs of $h$ --- i.e.\ the
\textbf{product} of factors, since $h$ is a log. Composite intermediates (screen
A against B, then the composite A$+$B against C) enter as additional pairs, so
the multiplicity matrix carries them with no special case
(\S\ref{sec:screening-pairs}).

% =====================================================================
\subsection{The tabular path}
\label{sec:tabular-path}

\code{compile.py:217-226} slices the flat weak table
(\S\ref{sec:weaktable}) into \code{(n\_rhoy, n\_temp)}, and the engine
interpolates with a clamped searchsorted:

\begin{srclisting}
i = maximum(0, minimum(len(rhoy)-1, searchsorted(rhoy, logrhoy)) - 1)
...
lam[col] = 10.0 ** _bilinear_vec(...)
\end{srclisting}

and raises on out-of-bounds rather than clipping --- a deliberate divergence
from MESA's behaviour (\S\ref{sec:weaktable}), chosen so the engine cannot
silently produce edge values.

\textbf{And the \rhoye{} that indexes it comes from the composition:}

\begin{srclisting}
ye = (cn.stoich.Z @ Y) / (cn.stoich.A @ Y)     # engine.py:151
logrhoy = np.log10(rho * ye)
\end{srclisting}

Note \Ye{} is computed as $\sum Z_iY_i / \sum A_iY_i$, i.e.\ normalised,
matching pynucastro's \code{Composition.ye}. Since $\sum A_iY_i = 1$ for a valid
composition this is a no-op --- but it is a \emph{robust} no-op, and it means
the tabular path degrades gracefully on slightly unnormalised inputs instead of
drifting.

% =====================================================================
\subsection{The lepton ledger as an assertion}
\label{sec:ledger-enforcement}

The ledger of \S\ref{sec:ledger} is enforced, not documented:

\begin{srclisting}
if dq != d_e[j]:
    raise ValueError("... nuclear dZ = ... but ledger d_electron = ... -- "
                     "charge-to-lepton closure broken")
\end{srclisting}

\begin{result}
\textbf{The ledgers are assigned from \code{weak\_type}, never back-derived from
$Z\cdot\nu$.} The docstring says why: back-deriving would make the
charge-to-lepton closure a tautology, and the conservation gate would pass
vacuously. This is the single best example in the repo of a test being designed
so it \emph{can} fail.
\end{result}

% =====================================================================
\subsection{Duplicate-link resolution}
\label{sec:duplicate-links}

Where both a REACLIB fit and a tabular rate cover the same transition
(\S\ref{sec:tablefamilies}), \code{build\_rate\_collection} drops the ReacLib
member:

\begin{srclisting}
for rate in group:
    if isinstance(rate, pyna.rates.TabularRate): survivors.append(rate)
    else: full.remove_rate(rate)
\end{srclisting}

% =====================================================================
\subsection{The code path, end to end}
\label{sec:codepath}

\begin{figure}[H]
\centering
\begin{tikzpicture}[
  font=\footnotesize,
  stage/.style={draw, rounded corners=2pt, line width=0.7pt, align=left,
                inner sep=5pt, text width=0.80\textwidth, fill=resultrule!5},
  lbl/.style={font=\scriptsize\itshape, inner sep=2pt},
  arr/.style={-{Latex[length=2mm]}, line width=0.6pt}]

\node[stage] (a) {\code{isotopes\_mesa80.yaml}};
\node[stage, below=7mm of a] (b)
  {\code{graph/network.py}: \code{ReacLibLibrary} $+$
   \code{TabularLibrary(ordering=PINNED)}\\
   \code{linking\_nuclei} $\to$ duplicate resolution (tabular wins) $\to$
   disposition drop\\
   \code{build\_stoich} $\to$ \code{weak\_mask}, \code{weak\_type},
   \code{d\_electron}\,/\,\code{d\_neutrino}\,/\,\code{d\_antineutrino}};
\node[stage, below=7mm of b] (c)
  {\code{fluxes/compile.py}: tabular rates $\to$ (\code{tab\_cols},
   \code{tab\_rhoy}, \code{tab\_temp}, \code{tab\_rate2d})\\
   strong rates $\to$ coefficient tensor (\S\ref{sec:coefftensor})};
\node[stage, below=7mm of c] (d)
  {\code{fluxes/engine.py}:
   \code{lam[tab\_cols] = 10**bilinear(log rho*Ye, log T)}\\
   $R = $ \code{prefactor} $\cdot \rho^{\code{densexp}} \cdot \prod Y \cdot$
   \code{lam}\\
   \code{dye\_weak = (Z nu)[weak] @ R[weak]}};

\draw[arr] (a) -- node[lbl,right]{\code{load\_isotope\_table}} (b);
\draw[arr] (b) -- (c);
\draw[arr] (c) -- (d);
\end{tikzpicture}
\caption{The weak-sector code path, from the isotope list to the \Ye{} signal.}
\label{fig:codepath}
\end{figure}

That last line is invariant \#3 made computable: the physical \Ye{} signal
carried by the weak columns, asserted nonzero by
\code{test\_flux\_engine.py::test\_dye\_weak\_nonzero}
(\S\ref{sec:oracles-db}).

% =====================================================================
\subsection*{Self-check --- the compiled evaluator}
\addcontentsline{toc}{subsection}{Self-check --- the compiled evaluator}
\label{sec:selfcheck-evaluator}

\begin{aside}
Every question in this block arrives from a source self-check whose material was
re-homed here: item 1 from the source's Part-I block, items 2--4 from the
Part-II block, items 5--6 from the Part-V block.
Appendix~\ref{app:concordance} records the mapping.
\end{aside}

\begin{selfcheck}
\item A reaction has \code{prefactor} $1/2$ and \code{dens\_exp} 1. What is its
      arity, and what are its reactants? Now do the same for \code{prefactor}
      $1/6$.
\item Write out \code{lam} for a chapter-8 rate (3 reactants $\to$ 1 product)
      with two sets, symbolically, all the way to $R_j$ including
      \code{prefactor} and $\rho$ powers.
\item Why does \code{set\_starts} exist rather than a loop over columns? Estimate
      the speedup for \msn{151} (1518 columns, ${\sim}1700$ sets, 4096 states).
\item \code{compile\_network} raises if \code{set\_owner} is not sorted. What
      downstream computation would silently produce wrong answers if it were
      not?
\item The engine raises on out-of-table states while MESA clips. Argue which is
      correct for \emph{this project}, and name a situation where the other
      choice would be right.
\item The engine reads only column 5 of the weak tables. Name one quantity the
      project needs that lives in another column, and say which invariant it
      blocks.
\end{selfcheck}
